%
% 2-Page Project Proposal
% IEEE two-column format
% Compile: pdflatex proposal.tex
%

\documentclass[journal]{IEEEtran}

\usepackage[pdftex]{graphicx}
\usepackage[cmex10]{amsmath}
\usepackage{amssymb}
\usepackage{booktabs}
\usepackage{url}

\hyphenation{multi-lingual mod-er-ation}

\begin{document}

\title{FairMod: Constrained Reinforcement Learning for\\
       Equitable Multilingual Content Moderation}

\author{Rishabh~J.\\
        \textit{[Institution Name]}\\
        rishabh12j@gmail.com}

\maketitle


%=== ABSTRACT ===================================================================
\begin{abstract}
Online communities increasingly span multiple languages, yet automated content
moderation systems remain largely monolingual and rely on blunt binary decisions.
This project proposes \textit{FairMod}, a reinforcement learning agent for
Discord moderation framed as a Constrained Markov Decision Process (CMDP).
The agent learns a five-level graduated enforcement policy---\textsc{Allow},
\textsc{Warn}, \textsc{Delete}, \textsc{Timeout}, \textsc{Ban}---across 13
languages, using a 403-dimensional observation space that combines multilingual
sentence embeddings, calibrated toxicity scores, user infraction history with
cool-down decay, and server-level heat signals. Procedural fairness is enforced
through hard action masking (bans require prior timeouts; timeouts require prior
infractions), and cross-language equity is maintained via an adaptive Lagrangian
multiplier and a Disparate Impact Wrapper that penalizes disproportionate ban
rates against any single language group. Evaluation on a held-out test split of
595 synthetic episodes shows 91.7\% accuracy on high-stakes enforcement actions
(Delete, Timeout, Ban)---categories where static baselines score 0--4.5\%---with
a false negative rate below $10^{-4}$. Per-language fairness is achieved for 12
of 14 language categories.
\end{abstract}

\begin{IEEEkeywords}
constrained MDP, content moderation, fairness, Lagrangian RL, multilingual NLP,
proximal policy optimization
\end{IEEEkeywords}


%=== I. PROBLEM & MOTIVATION ====================================================
\section{Problem Statement and Motivation}

\IEEEPARstart{A}{utomated} content moderation in multilingual communities faces
two unsolved challenges. First, enforcement should be \textit{graduated}: a
first-time mild offender deserves a warning, not a ban; a user who ignores three
warnings warrants escalation. Keyword filters and static toxicity thresholds
cannot model this because they are stateless and binary. Second, enforcement
should be \textit{equitable}: the same behavior should produce the same response
regardless of whether the user writes in English, Arabic, or Hinglish.
Neural toxicity classifiers trained primarily on English data produce
systematically lower scores for Arabic ($-0.139$ vs.\ English), Italian
($-0.096$), and Ukrainian ($-0.095$), and inflate scores for Hinglish
($+0.095$), translating directly into disparate moderation outcomes
\cite{conneau2020, dementieva2023}.

Reinforcement learning is a natural fit: the agent maintains state (user
history, server context) and learns a policy that maps the full observational
context to an enforcement action, optimizing for long-term community health
rather than per-message decisions. Framing the problem as a CMDP
\cite{altman1999} allows fairness constraints to be expressed as explicit
inequality constraints on the expected cost, solved via Lagrangian relaxation
\cite{stooke2020} without sacrificing primary task performance.


%=== II. PROPOSED SYSTEM ========================================================
\section{Proposed System}

\subsection{CMDP Formulation}

The moderation task is modeled as $\mathcal{M} = (S, A, P, R, C, d, \gamma)$
with $\gamma = 0.95$. The action space $A = \{0,\ldots,4\}$ maps to
\{\textsc{Allow}, \textsc{Warn}, \textsc{Delete}, \textsc{Timeout}, \textsc{Ban}\}.
The constrained objective is:
\begin{equation}
  \max_\pi\;\mathbb{E}_\pi\!\left[\sum_t \gamma^t R_t\right]
  \quad\text{s.t.}\quad
  \mathbb{E}_\pi\!\left[\sum_t \gamma^t C_t\right] \leq d
\end{equation}
where $C_t = \mathbf{1}[s_t < 0.30 \wedge a_t > 0]$ penalizes punitive actions
on benign content (false positives), and $d = 0.05$ is the maximum tolerable
false-positive rate.

\subsection{Observation Space}

The 403-dimensional observation vector is:
$\mathbf{o}_t = [\,\mathbf{e}_t \;\|\; s_t \;\|\; \mathbf{h}_t \;\|\;
\mathbf{q}_t \;\|\; \mathbf{l}_t\,]$,
where $\mathbf{e}_t \in \mathbb{R}^{384}$ is a MiniLM-L12-v2 sentence
embedding~\cite{reimers2019}, $s_t \in [0,1]$ is a per-language calibrated
XLM-R toxicity score, $\mathbf{h}_t \in \mathbb{R}^3$ encodes warn count,
timeout count, and effective infractions with cool-down decay,
$\mathbf{q}_t \in \mathbb{R}^2$ encodes server-level heat, and
$\mathbf{l}_t \in \{0,1\}^{13}$ is a language one-hot vector.

\subsection{Fairness Mechanisms}

Three complementary fairness layers are stacked:

\textbf{Action Masking.}
Hard constraints enforce procedural escalation: \textsc{Ban} is illegal until
at least one prior timeout exists; \textsc{Timeout} requires effective
infractions $\geq 2$; punitive actions are blocked when $s_t < 0.15$.

\textbf{Lagrangian Penalty.}
The reward is augmented as $\tilde{R}_t = R_t - \lambda C_t$, where $\lambda$
is updated online via dual ascent: $\lambda \leftarrow \max(\epsilon,\,
\lambda + \alpha_\lambda(\hat{C} - d))$.
This suppresses false positives across all languages uniformly.

\textbf{Disparate Impact Wrapper.}
A rolling window monitors per-language \textsc{Ban} rates. When language $l$
exceeds $1.5\times$ the overall rate, an additional $-2.0$ reward penalty is
applied, directly discouraging over-policing of specific language groups.

\subsection{Multilingual Toxicity Pipeline}

Raw XLM-R scores are corrected by a per-language affine calibration
$\hat{s} = \operatorname{clip}(s_{\text{raw}} \cdot \sigma_l + \beta_l, 0, 1)$
where scale $\sigma_l$ and offset $\beta_l$ are computed from percentile anchors
in the training corpus with English as the reference language. A compiled
multilingual regex module covering all 13 languages provides hard overrides for
explicit threats that the neural classifier under-detects; coverage is 13/13
languages for stalking and 13/13 for death threats after adding reverse
subject-object-verb pattern variants for Russian, Ukrainian, Spanish, and Hindi.

\subsection{Training Setup}

Training follows a two-stage pipeline inspired by~\cite{liu2025meta}: (1)~a behavior
cloning stage imitation-learns an oracle escalation rule via NLL minimization to anchor
the policy in the correct behavioral structure; (2)~an RL fine-tuning stage refines the
BC-initialized policy via MaskablePPO~\cite{huang2022mask} with a composite reward
combining per-step accuracy and a rubric-based escalation-consistency component.

The agent is trained on 336 synthetic Discord episodes (8,143 steps) generated by
prompting DeepSeek~V3.1 with real-data seeds from the Jigsaw and textdetox multilingual
corpora \cite{jigsaw, dementieva2023}. Thread types span five categories: toxic (30\%),
escalating (20\%), mild frustration (15\%), subtle (10\%), and benign (25\%).
Hinglish threads use a specialized code-switching prompt to ensure
Hindi--English mixed-script authenticity. Episodes are split 80/20 into train and test
sets, and checkpoint selection uses a constraint-penalized score that accounts for both
reward and false-positive rate.


%=== III. PRELIMINARY RESULTS ===================================================
\section{Preliminary Results}

Table~\ref{tab:results} summarizes evaluation on 100 held-out episodes (6,127
steps) using the current best checkpoint. The RL agent substantially outperforms
both baselines and is the only system capable of graduated enforcement.
All 13 language groups pass the $1.5\times$ ban-rate parity threshold in the
latest checkpoint, including Hinglish which previously exhibited a $2.60\times$
disparity before calibration correction and DisparateImpactWrapper activation.

\begin{table}[!t]
  \renewcommand{\arraystretch}{1.3}
  \caption{Comparative Evaluation Results}
  \label{tab:results}
  \centering
  \begin{tabular}{lcccc}
    \toprule
    System & Acc. & FP & FN & TIMEOUT/BAN \\
    \midrule
    Keyword Filter    & 67.3\% & 0.020 & 0.272 & No \\
    Static Threshold  & 77.0\% & 0.000 & 0.000 & No \\
    \textbf{FairMod}  & \textbf{73.6\%} & 0.153 & $\mathbf{10^{-4}}$ & \textbf{Yes} \\
    \bottomrule
  \end{tabular}
\end{table}


%=== IV. CURRENT STATUS & NEXT STEPS ============================================
\section{Current Status and Next Steps}

Since the initial proposal, the following have been implemented:
(1)~two-stage training (behavior cloning $\to$ RL) anchors the policy in the correct
escalation structure before refinement; (2)~a production \texttt{discord.py} bot with
persistent SQLite user state and a full audit trail enables live community testing;
(3)~Monte-Carlo confidence estimation via stochastic multi-sample voting improves
borderline decisions; (4)~wrapper ordering was corrected so that reward scaling is
applied \emph{after} all penalty terms; (5)~WARN action masking at $s_t < 0.15$ prevents
false-positive warnings on clearly benign content; and (6)~constraint-aware checkpoint
selection saves the best model based on reward \emph{and} false-positive rate, not reward
alone.

Remaining priorities are: (a)~expand threat coverage for mild-toxicity categories where
the XLM-R classifier produces near-zero signal in Italian and Arabic;
(b)~disagreement-based episode reweighting for improved data efficiency; and
(c)~extending the observation space with channel-graph context to distinguish
cross-channel escalation patterns.


%=== REFERENCES =================================================================
\begin{thebibliography}{9}

\bibitem{altman1999}
  E.~Altman, \textit{Constrained Markov Decision Processes}.
  CRC Press, 1999.

\bibitem{stooke2020}
  A.~Stooke, J.~Achiam, and P.~Abbeel,
  ``Responsive safety in reinforcement learning by PID Lagrangian methods,''
  in \textit{Proc.\ ICML}, 2020.

\bibitem{schulman2017}
  J.~Schulman \textit{et al.},
  ``Proximal policy optimization algorithms,''
  \textit{arXiv:1707.06347}, 2017.

\bibitem{huang2022mask}
  S.~Huang and S.~Onta\~n\'on,
  ``A closer look at invalid action masking in policy gradient algorithms,''
  in \textit{Proc.\ FLAIRS}, 2022.

\bibitem{conneau2020}
  A.~Conneau \textit{et al.},
  ``Unsupervised cross-lingual representation learning at scale,''
  in \textit{Proc.\ ACL}, 2020.

\bibitem{dementieva2023}
  D.~Dementieva \textit{et al.},
  ``Overview of the 1st shared task on multilingual text detoxification,''
  in \textit{Proc.\ ACL Workshop NLP4PI}, 2023.

\bibitem{reimers2019}
  N.~Reimers and I.~Gurevych,
  ``Sentence-BERT: Sentence embeddings using Siamese BERT-networks,''
  in \textit{Proc.\ EMNLP}, 2019.

\bibitem{jigsaw}
  Google Jigsaw,
  ``Toxic comment classification challenge,'' Kaggle, 2018.
  [Online]. Available: \url{https://kaggle.com/c/jigsaw-toxic-comment-classification-challenge}

\bibitem{liu2025meta}
  Z.~Liu \textit{et al.},
  ``Scaling reinforcement learning for content moderation with large language models,''
  \textit{arXiv:2512.20061}, 2025.

\end{thebibliography}

\end{document}
